\documentclass[a4paper,11pt]{article}
\usepackage{fancyhdr,graphics,graphicx,xy}
\usepackage{blkarray}
\usepackage{ifthen}
\usepackage{amssymb,amsmath,mathtools,color}
\pagestyle{fancy}
\usepackage{xparse}


%\textwidth = 6.5 in
%\textheight = 9 in
%\oddsidemargin = 0.0 in
%\evensidemargin = 0.0 in
%\topmargin = 0.0 in
%\headheight = 0.0 in
%\headsep = 0.0 in
%\parskip = 0.2in
%\parindent = 0.0in


%%%%%%%%%%%%%%%%%%%%%%%%%%%%%%%%%%%%%%%%%%
%%% General math macros
%%%%%%%%%%%%%%%%%%%%%%%%%%%%%%%%%%%%%%%%%%

\newcommand{\tf}{\tfrac} 
\renewcommand{\Re}{\mathrm{Re}}

\newcommand{\Def}{\coloneqq}
\newcommand{\lgg}{\operatorname{lg}}
\newcommand{\R}{\mathbb{R}}
\newcommand{\C}{\mathbb{C}}
\newcommand{\D}{\mathbb{D}}
\newcommand{\A}{\mathcal{A}}
\newcommand{\F}{\mathcal{F}}
\newcommand{\Q}{\mathbb{Q}}
\newcommand{\N}{\mathbb{N}} 
\newcommand{\Quat}{\mathbb{J}}
\newcommand{\Z}{\mathbb{Z}}
\newcommand{\T}{\mathbb{T}}
\newcommand{\U}{\mathcal{U}}
\newcommand{\LP}{\operatorname{L}}
\newcommand{\abs}[1]{\left| {#1} \right|}
\newcommand{\sgn}{\operatorname{sgn}}

\newcommand{\tendsto}{\to}
\newcommand{\tendto}{\to}

%%%%%%%%%%%%%%%%%%%%%%%%%%%%%%%%%%%%%%%%%%
%%% Probability Macros
%%%%%%%%%%%%%%%%%%%%%%%%%%%%%%%%%%%%%%%%%%
\newcommand{\stdcg}{N_\bbc(0,1)}

\newcommand{\expec}{\mathbb{E}}
\newcommand{\expect}{\mathbb{E}}
\newcommand{\Exp}{\mathbb{E}}
\newcommand{\E}{\mathbb{E}}
\newcommand{\prob}{\mathbb{P}}
\newcommand{\filt}{\mathscr{F}}
\newcommand{\Prob}{\prob}
\renewcommand{\Pr}{\prob}
\newcommand{\Ent}{\operatorname{Ent}}
\newcommand{\Var}{\operatorname{Var}}
\newcommand{\Cov}{\operatorname{Cov}}
\newcommand{\dtv}{d_{TV}}
\newcommand{\weakto}{\Rightarrow}
\newcommand{\lawof}{\mathscr{L}}
\newcommand{\Bernoulli}{\operatorname{Bernoulli}}
\newcommand{\Binomial}{\operatorname{Binom}}
\newcommand{\Poisson}{\operatorname{Poisson}}
\newcommand{\Normal}{\operatorname{Normal}}
\newcommand{\Unif}{\operatorname{Unif}}
\newcommand{\As}{\ensuremath{\text{a.s.}}}
\newcommand{\distequals}{\overset{\mathrm{d}}{=}}


\newcommand{\asstnumber}{3}
\newcommand{\assigntop}{
  \fancyhead[L]{}
  \fancyhead[C]{\Large Math 356 -- Fall 2023\\ Assignment \asstnumber}
  \fancyhead[R]{ }
}

\newtheorem{theorem}{Theorem}
\newtheorem{corollary}[theorem]{Corollary}
\newtheorem{definition}{Definition}
\newcommand{\p}[1]{\mathbb{P}(#1)}
\newcommand{\px}[1]{\mathbb{P}_x(#1)}
\newcommand{\py}[1]{\mathbb{P}_y(#1)}
\newcommand{\I}[1]{\mathbf{1}_{#1}}
\newcommand{\e}{\mathbb{E}}

\begin{document}
\assigntop

\section*{Question 1 }
\subsection*{(a)}

Notice that  a roll of a fair die, letting $X_i = \text{ see a 4 or not}$ where the range is $\{0,1\}$, then $\prob(X_i = 1) = \frac{1}{6}$ and $\prob(X_i = 0) = \frac{5}{6}$, then, letting $p = \frac{1}{6}$
$$\prob(X_i = 1) = p = 1 - \prob(X_i = 0)$$
we see that $X_i \sim Ber(p)$, then letting $S_n = X_1 + X_2 + \ldots + X_n$, $S_n \sim Bin(n,p)$, thus we can apply the law of large numbers for binomial distributions. We want to show that the probability of the frequency of 4s be at least $17 \%$ converges to 0 as $n \to \infty$. Observe that 
\begin{align*}
	\prob\left ( \frac{S_n}{n} \geq 0.17 \right) &= \prob\left(\frac{S_n}{n} - 0.1\bar{6} \geq 0.00\bar{3}\right)\\
	&= \prob\left(\frac{S_n}{n} - p  \geq 0.00\bar{3} \right)\\
	&= 1 - \prob \left(\frac{S_n}{n} - p < 0.00\bar{3}\right) \to 0 
	\end{align*}
	as $n$ tends to infinity, by the law of large numbers for binomial distributions. 
	\subsection*{(b)}
 Let $B_{n,i} \subseteq \mathcal{F}$ be the event such that the frequency of seeing a $i \in \{1,2,3,4,5,6\}$ is between $16\%$ and $17\%$, notice that $$A_n = \bigcap_{i=1}^6 B_{n,i}$$
 and
 $$A_n^c = \bigcup_{i=1}^6 B_{n,i}^c$$
 We want to show that for $n$ large enough
 $$ \prob \left( A_n^c \right) < 0.001$$ 
 While knowing that 
 $$ \prob(B_{n,i}^c) \to 0$$ 
 since the random variable associated with the number of any $i$s is the same as $S_n$, then since
 \begin{align*}  
     \prob(\tfrac{S_n}{n} \geq 0.16) \geq \prob(|\tfrac{S_n}{n} - p| \leq 0.00\bar{6} ) \to 1  \quad 
 \end{align*}
 \footnote{the event such that the distance between the frequency
 and $p$ is less than $0.00\bar{6}$ is a subset of the event that the frequency is greater than than $0.16$  }

 as $n$ goes to infinity thus the probability of the frequency of observing $i$s for any $i$ to be at least $16\%$ tends to $1$. From the converse of $a$, the probability of the frequency of observing $i$s for any $i$ to be at most $17\%$ tends to $1$, then 
 $$\prob(B_{i,n}^c) \to 0$$
 for all $i$. 
 Then it's straightforward to see that 
 $$\prob(A_n^c) = \prob\left(\bigcup_{i=1}^6 B_{n,i}^c \right) \leq \sum_{i=1}^6 \prob(B_{n,i}^c) \to 0$$
 for $n$ large enough, here we use countable (finite) subadditivity. 
	

\clearpage

\section{ Question 2 Solution}
Intuitively, the binomial distribution $S_n \sim Bin(n,p)$ can be approximated by a normal distribution with mean $np$ and variance $npq$. The peak of a normal distribution is achieved at it's mean. Since the probability mass function is discrete (we take it's value on the positive integers), we can show that it's first increasing then decreasing by investing the ratio between $P(S_n = k)$ and $P(S_n = k+1)$ given by
\begin{align*}
	\dfrac{n!}{\left( k+1\right) !\left( n-k-1\right) !}p^{k+1}q^{n-k-1} \cdot \dfrac{k!\left( n-k\right) !}{n!p^{k}q^{n-k}} = \frac{p(n-k)}{q(k+1)}
\end{align*}
$$
 \frac{p(n-k)}{q(k+1)}
$$
Which is $> 1$ when $k < p(n+1)-1$ and less than $1$ for $k > p(n+1)-1 $, which shows that the probability mass function is first increasing, then decreasing for $k$ greater than some amount depending on $n$ and $p$. Knowing this we know that the maximum probability is given by the value $k$ such that 
$$\frac{P(k)}{P(k-1)} > 1$$
	and
$$\frac{P(k+1)}{P(k)} < 1$$
two values $k$ and $k+1$ represent the same maximum when 
$$\frac{P(k+1)}{P(k)} = 1$$
In other words, the probability mass function attains a maximum at 
$$ \max_{k}(k < p(n+1)-1) + 1$$ 
if $p(n+1)-1 \notin \mathbb{Z}$, otherwise the probability mass function attains a maximum at 
$$ p(n+1)-1 \text{ and } p(n+1) \text { if the former is integral } $$

. The probability mass function of the poisson distribution, letting $X \sim Poi(\lambda)$ is given by
$$ P(X=k) = e^{-\lambda}\frac{\lambda^k}{k!} \quad k \in \{0,1,2,\ldots\}$$
The ratio for $P(X =k)$ and $P(X = k+1)$ is given by 
\begin{align*}
	\frac{P(X = k+1)}{P(X=k)} = e^{-\lambda}\frac{\lambda^{k+1}}{(k+1)!} \cdot \frac{k!}{e^{-\lambda}\lambda^k} = \frac{\lambda}{k+1}
\end{align*}
This ratio is $>1$ when $ k < \lambda - 1$ and $<1$ when $k > \lambda -1$, thus it's first increasing and then decreasing. The maximum is
$$\max_{k}(k < \lambda - 1)$$
if $\lambda -1 \in \mathbb{Z}$, otherwise, the maximum is attained at 
$$\lambda - 1 \text{ and } \lambda \text{ where the former is integral} $$

\clearpage


\section*{Question 3}
\subsection*{(a)}
\begin{align*}
	\expec(Y(Y-1)\cdots (Y-n+1)) &= \sum_{k=1}^\infty k(k-1)\cdots(k-n+1)e^{-\mu}\frac{\mu^k}{k!} \\
				     &= \sum_{k=n}^\infty e^{-\mu} \frac{\mu^k}{(k-n)!} \quad \text{since first $n-1$ terms vanish} \\
				     &= \mu^n \sum_{k=0}^\infty e^{-\mu}\frac{\mu^k}{k!} \\
				     &= \mu^n
\end{align*}
\subsection*{(b)}
We can calculate $\expec(Y^3)$ using linearity of expectation, observe that 
\begin{align*}
	\expec(Y^3) &= \expec(Y(Y-1)(Y-2)) + \expec(3Y^2) - \expec(2Y) \\
		    &= \mu^3 + 3(\mu^2 + \mu) - 2\mu \\
		    &= \mu^3 + 3\mu^2 + \mu
\end{align*}

\clearpage

\section*{Question 4}
This can be solved using a binomial distribution. Let $X_i$ be the random variable that represents if a persons birthday is the same as me or not. All the $X_i$ are independent and $\sim Ber(p)$ with $p = \frac{1}{365}$. Letting $X = X_1 + X_2 + \ldots X_i$ be the amount of people that have the same birthday as me, we know that $X \sim bin(n,p)$, we want to know $\mathbb{P}(X\geq 1) = 1- \prob(X=0)$
where
\begin{align*}
	1 - \prob(X = 0) &= 1 - \binom{n}{0} (1-\frac{1}{365})^n \\
			 &= 1 - (1-\frac{1}{365})^n \\
			 &\approx 1 - (1-\frac{1}{365})^n \\
			 &= 1 - e^{\frac{-n}{365}}
\end{align*}
Solving the inequality we get $1 - e^{\frac{-n}{365}} > \frac{2}{3} \iff \ln(\frac{1}{3}) > \frac{-n}{365} \iff 0 - \ln(3) > \frac{-n}{365} \iff n > 365\ln(n) \sim 401$
\clearpage


\section*{Question 5}
Jessica flipping 5 times a faircoin per day, letting $X_i$ be the random variable such that it's equal to $1$ when all flips comes up tails, and $0$ otherwise, then $X_i \sim Ber(p)$ with $p = \frac{1}{2}^5 = \frac{1}{32}$, and $n = 30$. Observe that $np^2 = 30 (\frac{1}{32}^2) \approx 0.03 $ thus the poisson distribution approximates the binomial pretty well (theorem 4.20 of the textbook), but observe that $np(1-p) = 30 \cdot \frac{1}{32} \cdot {31}{32} \approx 1$ is far below 10 (textbook), so the Poisson distribution should work better here.  Therefore, we use the poisson distribution, where $S_n  = X_1 + X_2 + \ldots + X_n \sim Bin(n,p)$ is approximated by, letting $\lambda = np = \frac{30}{32}$, $e^{-\lambda} \frac{\lambda^k}{k!}$
then 
$$ \prob(S_n = k) = \frac{e^{-\frac{30}{32}} \left(\frac{30}{32}\right)^2}{2} \approx 0.172$$


\clearpage

\section*{Question 6}

Observe that each $F_n$ can be thought of as a step function where $F_n(t) = \frac{1}{n}$ for $t \in [0,1)$, and $F_n(t) = \frac{2}{n}$ for $t \in [1,2)$ and so on, where $F_n(t) = \frac{k+1}{n}$ for $t \in [k,k+1)$, that means 
$$F(t) = \frac{\lfloor t \rfloor + 1}{n}$$
it's clear that the limit of $F(t)$ as $n$ goes to infinity is $0$, 
\end{document}

